\noindent
This chapter indexes the 30 catalogue records assigned to 1974--75 in the supplied
package, with 265 mathematical occurrences, 10 definite problem entries,
78 bibliography entries and 92 critical corrections or qualifications.
Records retain their catalogue order; items follow the surviving text within
each record. Status statements refer to November 2026. The classifications
describe historical fate; the separate assessments describe literal correctness,
publication matching and verification.

\begin{center}
\begin{tabular}{@{}lrrrrrr@{}}
\toprule
Class & A & B & C & D & E & F\\
\midrule
Occurrences & 29 & 10 & 0 & 0 & 203 & 23\\
\bottomrule
\end{tabular}
\end{center}

\paragraph{The main mathematical lines.}
The finite-gap sequence connects Floquet theory, Riccati densities, stationary
KdV equations, hyperelliptic curves and Jacobians. Its chronology matters:
Novikov's paper was received in February and published in July 1974, before
the September meetings. The joint Dubrovin--Novikov announcement, received in
June and published later in 1974, gives the converse and algebraic-geometric
synthesis. The later survey supplies fuller hypotheses, not the first
circulation of all its results
\citep{novikov1974periodic,dubrovinnovikov1974periodicity,dubrovinmatveevnovikov1976nonlinear}.
Kushnirenko's local Milnor-number and global critical-point questions have
different inequality directions and nondegeneracy conditions
\citep[Theorems I--II]{kushnirenko1975newton}.

Giusti's lectures juxtapose scalar regularity, vectorial counterexamples and
relaxed variational problems. Segal moves from configurations and cancellation
to nerves, buildings and classifying spaces; the categorical material continues
1973-308 and 1973-309. Graev's main construction explains work already published
in 1973. Tulovskii's sequence is anchored in the 1971--72 FIO theory
\citep{hormander1971fourier,duistermaathormander1972fourierii}.
Sinai's phase-diagram lecture has a close subsequent Pirogov--Sinai publication,
whereas Dobrushin's conditional-law formulation predates the talk
\citep{pirogovsinai1975phase,dobrushin1968conditional}.
King's lectures explain divisor currents, counting functions and curvature
inequalities in relation to existing holomorphic-mapping theory
\citep{griffithsking1973nevanlinna,griffiths1971canonical}.

\paragraph{What the critical checks change.}
The notebook is often correct shorthand, but hypotheses and conventions cannot
be discarded. On P000555 the Hamiltonian quartic coefficient needs \(5/8\);
the copied \(5/2\) gives a nonzero Poisson bracket. Its spectral quintic has
a further normalization problem. Giusti's copied tensor on P000565 fails
the claimed radial-solution substitution. The two logarithmic integrals on
P000564 behave differently: the first has the wrong sign, while the symmetric
golden-ratio integral is correct. The Steinberg module requires reduced
homology at \(n=2\). Fourier signs, radial characteristics, Haar measures and
the divisor-counting boundary term provide further convention-sensitive checks.

\paragraph{Questions and the scope of development.}
No item is assigned C or D, and no Stage 2 research problem meets the selection
rule. This does not imply that the surrounding subjects contain no open problems.
Definite questions have established answers in specified settings; several
prompts lack enough data for an exact formulation. Entire functions of order
one, an elliptic analogue of resolution and the \(p\)-adic analogies are kept
separate from later thematic work. The verification section gives complete
local arguments, counterexamples to overbroad wording and independently checked
computations. These are editorial audits of old mathematics, without novelty claims.

\paragraph{Source boundaries and earlier years.}
P000551 is headed 25 March 1974, although records 331--332 are catalogued as
25 March 1975. Their five entries duplicate 1973-331.1--4 and 1973-332.1.
The requested identifiers and catalogue dates are retained as compatibility
aliases with explicit warnings. Removing these exact source duplicates leaves
260 occurrences newly indexed by this package. P000550 belongs to the preceding
finite-gap talk; P000552--553 belong to the 1 April 1974 spectral lecture.
Their earlier-index coverage is audited separately. Further links include
1967-132.7 on the ray space, 1973-326 on finite-gap inverse data and monotonicity,
and the current-group questions of 1969. Related problems are not merged when
the earlier index supplies only titles and cannot establish identical hypotheses.
